% ---
% id: reductions/rem:channel-choice
% kind: rem
% title: "Soundness relative to the channel, fixed by the carried tuple"
% summary: >-
%   Channels realizing one squashing relation can have different reachable
%   sets, so constraints sound for different channels may not be combined. A
%   constraint built from an operator in the carried tuple has the same
%   pull-back for every realizing channel.
% prominence: standard
% status: accepted
% uses: [lem:statistics, thm:transfer, cor:sound, def:squashing, cor:transported]
% mentions: [req:carried-tuple-emitted]
% assumes: []
% refs: []
% provenance: newly-derived
% lean:
% ---

\begin{remark}\label{rem:channel-choice}
The relation $M\dsq F$ asserts that \emph{some} squashing channel exists, and
different channels realizing it have different reachable sets: by
Lemma~\ref{lem:statistics}, two channels realizing
$(M;O)\dsq(F;\hat O)$ agree only on
$V\coloneqq\operatorname{span}\big(\{F_v\}\cup\{\hat O_k\}\big)$, and nothing
forces $\Srch_E(\Lambda)=\Srch_E(\Lambda')$. Three consequences follow.
\emph{(a)} Theorem~\ref{thm:transfer} needs only one channel, and every
realizing channel gives a valid conclusion, since $\Lambda$ does not appear in
it. \emph{(b)} There is something to choose between them for
\emph{tightness}: the reachable sets are in general incomparable, so which
channel is used decides which constraints are available. \emph{(c)}
Corollary~\ref{cor:sound}\,(ii)--(iii) therefore require a \emph{common}
channel: constraints sound for $\Lambda$ and constraints sound for $\Lambda'$
may have an intersection excluding every state reachable under either, and
imposing them together would void Theorem~\ref{thm:transfer} exactly as an
unsound constraint does.

What removes the hazard is the tuple of \tref{carried operators} of
Definition~\ref{def:squashing}. If $\hat O^{(n)}$ lies in the span of the
$n$-fold products of elements of $V$, then
$(\Lambda^{\dagger})^{\otimes n}(\hat O^{(n)})
=(\Lambda'^{\dagger})^{\otimes n}(\hat O^{(n)})$, since the adjoints act
factorwise and agree on $V$; the pull-back, and with it every constraint
Corollary~\ref{cor:transported} builds from it, is then the same for every
realizing channel. This is what the tuple is for: listing an operator in the
relation promotes a constraint that would otherwise be tied to one construction
into one that belongs to the relation itself, which is why the emitted
constraints are built from the tuple
(Requirement~\ref{req:carried-tuple-emitted}).
\end{remark}
